% GENERATED FILE. Edit canonical lessons, then run npm run generate:books.
% canonical-source-hash: fnv1a64:b66e378e8bafd4a0
% canonical-lessons: ZH-C473-zui, ZH-C473-yayi, ZH-W473-ya-tooth, ZH-C473-yachi, ZH-C473-baya, ZH-C473-buya

\chapter{At the Dentist}
\label{ch:zh-at-the-dentist}
\hlchaptermodality{473}

\begin{chapteropening}
\textbf{By the end of this chapter:} \emph{I can talk about my mouth, my dentist and my teeth, and having a tooth out or filled.}

Complete the last lesson of chapter 473: I can talk about my mouth, my dentist and my teeth, and having a tooth out or filled.
\end{chapteropening}

\section[zuǐ]{zuǐ — mouth; beak}

\label{lesson:ZH-C473-zui}



\begin{quote}
Before the new one: say the Mandarin for the look on one's face, then the Mandarin for the forehead.
\end{quote}

\subsection*{You'll want to know: zuǐ}
\emph{zuǐ} — \textquotedblleft{}mouth; beak\textquotedblright{}.

\begin{quote}
\emph{zhège cài tài là le wǒ de zuǐ hěn nánshòu} — This dish is so hot my mouth is burning.
\end{quote}

Its characters come later; for now it is a word for the ear.

\emph{zuǐ} is the everyday word for the mouth; \zh{口} \emph{kǒu} mostly appears inside other words, as in \emph{rùkǒu} 'entrance'.

\begin{cousinweb}
\emph{zuǐ} No dictionary we use records an older form, so its history is not settled.
\end{cousinweb}

\begin{grammarlens}[title={What you've built}]
One more word for saying how it was.
\end{grammarlens}

\subsection*{Your turn}
\begin{itemize}
\raggedright
  \item \emph{Say it:} \emph{zuǐ}
  \item \emph{Say it:} \emph{zuǐ}, once more
  \item \emph{Say it:} \emph{zhège cài tài là le wǒ de zuǐ hěn nánshòu} — This dish is so hot my mouth is burning.
  \item \emph{Recall:} say the Mandarin for the look on one's face, then the Mandarin for the forehead, then say \emph{zuǐ} again
\end{itemize}

\begin{tcolorbox}[breakable,colback=teal!4,colframe=teal!35!black,title={Before you move on}]
What does zuǐ mean? (\textquotedblleft{}Mouth; beak\textquotedblright{}.) Say it once more.
\end{tcolorbox}
\section[\texorpdfstring{\zh{牙医} (yáyī)}{牙医 (yáyī)}]{\zh{牙医} (yáyī) — dentist}

\label{lesson:ZH-C473-yayi}



\begin{quote}
Before the new one: say the Mandarin for the forehead, then the Mandarin for the mouth.
\end{quote}

\subsection*{You'll want to know: \zh{牙医}}
\textbf{\zh{牙医}} — \emph{yáyī} — \textquotedblleft{}dentist\textquotedblright{}.

\begin{quote}
\textbf{\zh{我的牙医很和气}} \emph{wǒ de yáyī hěn héqi} — My dentist is very nice.
\end{quote}

The new character, in pieces: \zh{牙} — four strokes: a short bar on top; a short stroke down on the left that turns and runs right; a vertical that hooks to the left at the foot; and a long stroke falling to the left. The next lesson puts it in your hand.

\begin{cousinweb}
\textbf{\zh{牙医}} (\emph{yáyī}) is built from \zh{牙} 'tooth' and \zh{医} 'doctor', both inherited from Middle Chinese, \emph{ngae} and \emph{'i} (Guangyun, in Baxter's transcription), which glosses the first 'teeth' and the second 'to heal'.
\end{cousinweb}

\begin{grammarlens}[title={What you've built}]
One more word for saying how it was.
\end{grammarlens}

\subsection*{Your turn}
\begin{itemize}
\raggedright
  \item \emph{Say it:} \emph{yáyī}
  \item \emph{Say it:} \emph{yáyī}, once more
  \item \emph{Say it:} \emph{wǒ de yáyī hěn héqi} — My dentist is very nice.
  \item \emph{Recall:} say the Mandarin for the forehead, then the Mandarin for the mouth, then say \emph{yáyī} again
\end{itemize}

\begin{tcolorbox}[breakable,colback=teal!4,colframe=teal!35!black,title={Before you move on}]
What does \zh{牙医} mean? (\textquotedblleft{}Dentist\textquotedblright{}.) Say it once more.
\end{tcolorbox}
\section[\texorpdfstring{\zh{牙} (yá)}{牙 (yá)}]{\zh{牙} — four strokes, each one different}

\label{lesson:ZH-W473-ya-tooth}



\begin{quote}
The word first, then the character inside it.

\begin{itemize}
\raggedright
  \item \emph{Recall:} say \emph{yáyī} — \textbf{R1}, one lesson back, before it decays
\end{itemize}

\begin{quote}
\zh{医} \zh{口}
\end{quote}

\emph{Read it:} each one aloud, with nothing said first — \emph{yī}, \emph{kǒu}
\end{quote}

\subsection*{Script — four strokes, each one different}
\hlblockfigure{\detokenize{figures/ZH-W473-ya-tooth-filmstrip.pdf}}{How \zh{牙} is written, stroke by stroke}

\begin{quote}
\zh{牙}
\end{quote}

\textbf{One piece of four strokes.} Count them as you go.

\textbf{The first two.} A short bar along the top, from left to right. Then under it on the left a short stroke down that turns and runs right, longer than the bar.

\textbf{The last two.} A vertical from the top bar down through the second stroke, hooking to the left at the foot. Then a long stroke falling to the left, from the middle down to the bottom left.

2 + 2 = four.

\subsection*{Your turn}
\begin{itemize}
\raggedright
  \item \emph{Write it:} \zh{牙}, all four, with the model in front of you
  \item \emph{Cover:} write it again from memory, and count the pieces, not the strokes
\end{itemize}

\begin{tcolorbox}[breakable,colback=teal!4,colframe=teal!35!black,title={Before you move on}]
How many strokes does \zh{牙} have? (\textbf{Four}.) Which comes last? (\textbf{The long stroke falling to the left}.)

Source: \href{https://raw.githubusercontent.com/chanind/hanzi-writer-data/68d10a4b21150cae5e1ebbd223eed289cf32d90c/data/\%E7\%89\%99.json}{Hanzi Writer Data \zh{牙}, pinned stroke paths and medians}.
\end{tcolorbox}
\section[yáchǐ]{yáchǐ — tooth; teeth}

\label{lesson:ZH-C473-yachi}



\begin{quote}
Before the new one: say the Mandarin for the mouth, then the Mandarin for a dentist.
\end{quote}

\subsection*{You'll want to know: yáchǐ}
\emph{yáchǐ} — \textquotedblleft{}tooth; teeth\textquotedblright{}.

\begin{quote}
\emph{tā de yáchǐ bǐ wǒ de bái} — His teeth are whiter than mine.
\end{quote}

Its characters come later; for now it is a word for the ear.

\begin{cousinweb}
\emph{yáchǐ} is built from \zh{牙} and \emph{chǐ}, both inherited from Middle Chinese, \emph{ngae} and \emph{tsyhiX} (Guangyun, in Baxter's transcription), which glosses each with this very word: 'teeth'.
\end{cousinweb}

\begin{grammarlens}[title={What you've built}]
One more word for saying how it was.
\end{grammarlens}

\subsection*{Your turn}
\begin{itemize}
\raggedright
  \item \emph{Say it:} \emph{yáchǐ}
  \item \emph{Say it:} \emph{yáchǐ}, once more
  \item \emph{Say it:} \emph{tā de yáchǐ bǐ wǒ de bái} — His teeth are whiter than mine.
  \item \emph{Recall:} say the Mandarin for the mouth, then the Mandarin for a dentist, then say \emph{yáchǐ} again
\end{itemize}

\begin{tcolorbox}[breakable,colback=teal!4,colframe=teal!35!black,title={Before you move on}]
What does yáchǐ mean? (\textquotedblleft{}Tooth; teeth\textquotedblright{}.) Say it once more.
\end{tcolorbox}
\section[báyá]{báyá — to have a tooth pulled}

\label{lesson:ZH-C473-baya}



\begin{quote}
Before the new one: say the Mandarin for a dentist, then the Mandarin for the teeth.
\end{quote}

\subsection*{You'll want to know: báyá}
\emph{báyá} — \textquotedblleft{}to have a tooth pulled\textquotedblright{}.

\begin{quote}
\emph{wǒ jīntiān qù báyá le xiànzài hái hěn téng} — I had a tooth out today, and it still really hurts.
\end{quote}

Its characters come later; for now it is a word for the ear.

\begin{cousinweb}
\emph{báyá} is built from \emph{bá} 'to pull up' and \zh{牙} 'tooth', both inherited from Middle Chinese, \emph{beat} and \emph{ngae} (Guangyun, in Baxter's transcription). The dictionary lists the first under three readings and glosses this one 'to pull up' and 'to promote'.
\end{cousinweb}

\begin{grammarlens}[title={What you've built}]
One more word for saying how it was.
\end{grammarlens}

\subsection*{Your turn}
\begin{itemize}
\raggedright
  \item \emph{Say it:} \emph{báyá}
  \item \emph{Say it:} \emph{báyá}, once more
  \item \emph{Say it:} \emph{wǒ jīntiān qù báyá le xiànzài hái hěn téng} — I had a tooth out today, and it still really hurts.
  \item \emph{Recall:} say the Mandarin for a dentist, then the Mandarin for the teeth, then say \emph{báyá} again
\end{itemize}

\begin{tcolorbox}[breakable,colback=teal!4,colframe=teal!35!black,title={Before you move on}]
What does báyá mean? (\textquotedblleft{}To have a tooth pulled\textquotedblright{}.) Say it once more.
\end{tcolorbox}
\section[bǔyá]{bǔyá — to have a tooth filled}

\label{lesson:ZH-C473-buya}



\begin{quote}
Before the new one: say the Mandarin for the teeth, then the Mandarin for to have a tooth pulled.
\end{quote}

\subsection*{You'll want to know: bǔyá}
\emph{bǔyá} — \textquotedblleft{}to have a tooth filled\textquotedblright{}.

\begin{quote}
\emph{wǒ zuì pà bǔyá} — Having a filling is what I dread most.
\end{quote}

Its characters come later; for now it is a word for the ear.

\begin{cousinweb}
\emph{bǔyá} is built from \emph{bǔ} 'to patch' and \zh{牙} 'tooth', both inherited from Middle Chinese, \emph{puX} and \emph{ngae} (Guangyun, in Baxter's transcription), which glosses the first 'to mend' and, after an older book, mending clothes.
\end{cousinweb}

\begin{grammarlens}[title={What you've built}]
One more word for saying how it was.
\end{grammarlens}

\subsection*{Your turn}
\begin{itemize}
\raggedright
  \item \emph{Say it:} \emph{bǔyá}
  \item \emph{Say it:} \emph{bǔyá}, once more
  \item \emph{Say it:} \emph{wǒ zuì pà bǔyá} — Having a filling is what I dread most.
  \item \emph{Recall:} say the Mandarin for the teeth, then the Mandarin for to have a tooth pulled, then say \emph{bǔyá} again
\end{itemize}

\begin{tcolorbox}[breakable,colback=teal!4,colframe=teal!35!black,title={Before you move on}]
What does bǔyá mean? (\textquotedblleft{}To have a tooth filled\textquotedblright{}.) Say it once more.
\end{tcolorbox}
